\subsection{Hard sync 的确定性误差与 teacher 漂移}
\label{sec:hs-noise}

固定 prefix，并允许 teacher 随 outer 轮次改变，记 $t_k=\pi_{T,k}(\cdot\mid x)$。把实际输出写成
\begin{equation}\label{eq:hs-noisy-update}
 \pi_{k+1}(a)\propto\pi_k(a)^{1-\eta_k}t_k(a)^{\eta_k}
             \exp\{e_k(a)\}.
\end{equation}
这里 $e_k$ 是输出相对精确 Gibbs target 的未归一化 log-density 误差。它可以包含有限 inner、目标计算和参数共享带来的偏差；本节先直接分析输出误差，不把梯度方差直接当作 $e_k$ 的方差。记 $\operatorname{osc}f=\max_a f(a)-\min_a f(a)$，并定义
\[
 b_k=d_H(\pi_k,t_k),\qquad E_k=\operatorname{osc}e_k,
 \qquad D_k=d_H(t_k,t_{k+1}).
\]
其中 $d_H$ 是前节定义的 log-ratio 跨度。$D_k$ 描述 teacher 分布漂移，与最后一节的 prefix 访问测度漂移分开。

\begin{theorem}[输出误差的输入--状态界]\label{thm:hs-input-state}
若所有分布严格为正，则式~\eqref{eq:hs-noisy-update} 满足
\begin{equation}\label{eq:hs-drifting-recursion}
 b_{k+1}\leq |1-\eta_k|b_k+E_k+D_k.
\end{equation}
固定 $0<\eta<2$，令 $c=|1-\eta|<1$，则
\[
 b_K\leq c^K b_0+\sum_{j=0}^{K-1}c^{K-1-j}(E_j+D_j).
\]
若 $E_j+D_j\leq B$，则 $\limsup_K b_K\leq B/(1-c)$。若 $E_j+D_j\to0$，则 $b_K\to0$。
\end{theorem}

\begin{proof}
对更新取相对 $t_{k+1}$ 的 log-ratio：
\[
 \log\frac{\pi_{k+1}}{t_{k+1}}
 =(1-\eta_k)\log\frac{\pi_k}{t_k}+e_k
       +\log\frac{t_k}{t_{k+1}}-\log Z_k.
\]
动作无关的 $\log Z_k$ 在跨度中抵消。由 $\operatorname{osc}(f+g)\leq\operatorname{osc}f+\operatorname{osc}g$ 和 $\operatorname{osc}(af)=|a|\operatorname{osc}f$，得到式~\eqref{eq:hs-drifting-recursion}。逐次代入递推得到卷积式，再利用 $\sum_{i=0}^{K-1}c^i=(1-c^K)/(1-c)$ 得到稳态上界。

对消失扰动，给定 $\delta>0$，选 $J$ 使 $j\geq J$ 时 $E_j+D_j\leq(1-c)\delta$。有限前缀贡献 $c^Kb_0+\sum_{j<J}c^{K-1-j}(E_j+D_j)$ 趋于零，剩余贡献至多 $\delta$。令 $\delta\downarrow0$ 得到收敛。
\end{proof}

该界解释了 hard sync 对误差的放大系数。$0<\eta\leq1$ 时分母为 $\eta$，$1\leq\eta<2$ 时分母为 $2-\eta$。相同的逐轮输出跨度误差在 $\eta=1$ 时具有最小最坏情形放大系数；接近稳定区间两端时误差积累更强。

固定 teacher、固定误差向量 $e$ 时，不动点可以精确求出：
\[
 \pi_\infty(a)=\frac{t(a)\exp\{e(a)/\eta\}}
 {\sum_v t(v)\exp\{e(v)/\eta\}},\qquad
 d_H(\pi_\infty,t)=\frac{\operatorname{osc}e}{\eta}.
\]
证明是在不动点更新中把 $\eta\log(\pi_\infty/t)=e-\text{常数}$ 解出。对于 $\eta\leq1$，它达到上述最坏情形稳态界。对于 $\eta>1$，固定误差的稳态比一般界小，因为一般界还允许误差随轮次交替；在居中 log-ratio 坐标中取 $e_k=(-1)^{k+1}e$，便可产生 $\operatorname{osc}e/(2-\eta)$ 的交替稳态跨度。

\subsection{随机输出误差的精确矩递推}

固定 teacher $t=\pi_T$，定义相对 teacher 的居中坐标
\[
 y_k(a)=\log\frac{\pi_k(a)}{t(a)}
       -\E_t\log\frac{\pi_k}{t},\qquad
 \xi_k(a)=e_k(a)-\E_te_k.
\]
$\E_ty_k=\E_t\xi_k=0$。居中操作是线性的，因此归一化常数完全消去，式~\eqref{eq:hs-noisy-update} 给出精确线性系统
\begin{equation}\label{eq:hs-centered-linear}
 y_{k+1}=(1-\eta)y_k+\xi_k.
\end{equation}
这个等式在大误差时也成立；把 $\|y_k\|^2$ 换成 KL 时才需要局部近似。

\begin{theorem}[Hard sync 的偏差与协方差传播]\label{thm:hs-covariance}
令 $a_\eta=1-\eta$。若 $y_0$ 确定、$\E[\xi_k\mid\mathcal F_k]=0$、$\E[\xi_k\xi_k^\top\mid\mathcal F_k]=\Sigma_k$ 为确定矩阵，且 $y_k$ 是 $\mathcal F_k$ 可测的，则
\[
 \E y_K=a_\eta^K y_0,\qquad
 \operatorname{Cov}(y_K)=\sum_{j=0}^{K-1}a_\eta^{2(K-1-j)}\Sigma_j.
\]
若 $0<\eta<2$ 且 $\Sigma_j=\Sigma$，则
\begin{equation}\label{eq:hs-stationary-cov}
 \operatorname{Cov}(y_K)\longrightarrow
       \frac{\Sigma}{\eta(2-\eta)}.
\end{equation}
若 $\xi_k=m+\nu_k$，其中 $\nu_k$ 条件均值为零，则 $\E y_K\to m/\eta$，协方差结论对 $\nu_k$ 保持成立。
\end{theorem}

\begin{proof}
展开式~\eqref{eq:hs-centered-linear}：
\[
 y_K=a_\eta^Ky_0+\sum_{j=0}^{K-1}a_\eta^{K-1-j}\xi_j.
\]
条件均值为零推出 $\E\xi_j=0$。对 $i<j$，$\xi_i$ 已是 $\mathcal F_j$ 可测，所以
\[
 \E[\xi_i\xi_j^\top]
 =\E[\xi_i\E(\xi_j^\top\mid\mathcal F_j)]=0.
\]
双重求和中的交叉项消失，只留下所给协方差和。常协方差时，平方几何级数的分母为 $1-a_\eta^2=\eta(2-\eta)$。非零均值情形直接对 $\E y_{k+1}=a_\eta\E y_k+m$ 求解。
\end{proof}

令 $M_t=\operatorname{diag}(t)$。由 $\pi_k=t\exp(y_k)/\E_t\exp(y_k)$，在 $\|y_k\|_\infty\to0$ 时有
\[
 \KL(\pi_k\Vert t)=\frac12 y_k^\top M_t y_k
                   +O(\|y_k\|_\infty^3).
\]
因此有限时间的二阶预期 KL 为
\[
 \frac12\left[(\E y_K)^\top M_t(\E y_K)
             +\operatorname{tr}\{M_t\operatorname{Cov}(y_K)\}\right].
\]
固定 $0<\eta<2$，若 $\nu_k$ 是独立高斯向量，并令 $m=\epsilon\bar m$、$\Sigma=\epsilon^2\bar\Sigma$ 随同一尺度 $\epsilon\downarrow0$ 缩小，其中 $\bar m,\bar\Sigma$ 固定，则稳态三阶绝对矩为 $O(\epsilon^3)$。上式的稳态二阶项就是
\begin{equation}\label{eq:hs-noise-kl-floor}
 \frac{1}{2\eta^2}m^\top M_tm
 +\frac{\operatorname{tr}(M_t\Sigma)}{2\eta(2-\eta)}.
\end{equation}
稳态预期 KL 等于该二阶项加 $O(\epsilon^3)$。式~\eqref{eq:hs-stationary-cov} 则是精确的 log-ratio 协方差结论。仅有二阶矩假设时不能宣称 KL 的高阶余项也受控。

误差注入的位置决定步长依赖。若每轮实际输出的误差 $\xi_k$ 具有固定协方差，式~\eqref{eq:hs-noise-kl-floor} 的无偏噪声底与 $[\eta(2-\eta)]^{-1}$ 成正比。若误差来自 reward 估计，并以 $\xi_k=\eta\zeta_k$ 进入指数，固定的是 $\zeta_k$ 的协方差 $\Sigma_\zeta$，则噪声底改为
\[
 \frac{\eta}{2(2-\eta)}\operatorname{tr}(M_t\Sigma_\zeta).
\]
第二种模型中减小步长降低稳态噪声，同时减慢确定性收敛。第三章的 KL estimator 无偏性控制相应 reward 或梯度误差的条件均值；从它到本节的 $m,\Sigma$ 还要经过 inner 优化和停止梯度规则，不能直接等同。

\subsection{同步操作与分布空间模型的对应关系}

\subsubsection{保留快速 outer 收敛所需的 inner 精度}

固定 teacher，记每轮相对正确 target 的输出跨度误差为 $E_k$。Gibbs 在 $c=|1-\eta|\in(0,1)$ 时满足 $b_{k+1}\le cb_k+E_k$。若 $E_k\le B\gamma^k$，其中 $0<\gamma<1$，则卷积可精确求和：
\[
 b_K\le c^Kb_0+
 \begin{cases}
 B(c^K-\gamma^K)/(c-\gamma),&c\ne\gamma,\\
 BKc^{K-1},&c=\gamma.
 \end{cases}
\]
所以输出误差衰减得比理想动力学慢时，outer 收敛率由 $\gamma$ 决定；相同衰减率时还会出现 $K$ 因子。若一个 inner 求解器能保证 $E_k\le C\beta^{N_k}$，其中 $0<\beta<1$，要达到 $E_k\le B\gamma^k$，充分的 inner 步数为
\[
 N_k\ge\max\left\{0,
 \frac{\log(C/B)+k\log(1/\gamma)}{\log(1/\beta)}\right\}.
\]
这里指数收敛的 inner 假设必须另行证明；前面的普通非凸 SGD 梯度范数界不提供这个保证。

Mean 在 $\eta=1$ 时，令 $H(b)=e^bb^2/(8+e^bb^2)$，三次跨度定理给出
\[
 b_{k+1}\le H(b_k)b_k+E_k.
\]
当 $b_k\to0$ 时，保留三次阶数需要 $E_k=O(b_k^3)$；若希望保留精确 Mean 的三次首项，则需要更强的 $E_k=o(b_k^3)$。结合前述最小 anchor 概率下的 excess-loss 转换，$E_k\le2\sqrt{2\delta_k/m_0}$，在所考察轮次具有统一 $m_0>0$ 时，一个充分条件是 $\delta_k=O(b_k^6)$，保留首项则取 $o(b_k^6)$。若只能给出逐轮下界 $m_{0,k}$，相应条件改为 $\delta_k=O(m_{0,k}b_k^6)$ 或 $o(m_{0,k}b_k^6)$。即使 $\eta=1$ 的精确 outer 极快，固定 inner 精度也会使这段快速收敛提前停止。

固定误差上界 $E_k\le E$ 时可以得到比一般 Mean 不变跨度界更精细的局部结论。若 $H(2E)\le1/2$ 且 $b_0\le2E$，由于 $H$ 非减，
\[
 b_{k+1}\le H(2E)\,2E+E\le2E.
\]
归纳得到不变邻域。对 $L=\limsup b_k$，递推进一步给出
\[
 L\le H(2E)L+E,
 \qquad L\le\frac{E}{1-H(2E)}=E+O(E^3)
 \quad(E\downarrow0).
\]
这个结论要求初值位于上述局部邻域；它没有替代前面适用于大跨度的全局扰动界。

\subsubsection{同步边界的冻结量}

本章 hard sync 表示把下一轮 anchor 完全设为本轮 inner 输出。一次参数复制满足这一定义，却不保证 inner 输出已等于逐 prefix 的总体最优分布。共享参数、梯度裁剪、有限 batch 与有限 inner 都进入 $e_k$；若梯度采样规则改变了总体目标，$e_k$ 还包含持续的目标偏差。

同时冻结 rollout 采样测度、anchor log-probability、reward 和其 KL center，才形成前两章的固定 inner 子问题。只更新 anchor 参数而保留上一轮的 log-probability 或 center，会形成不一致的子问题。下一轮必须在同一同步边界上更新这些量，或者把过期误差显式计入本章的递推。

分布空间的可解性与参数空间的可实现性分开检查。逐 prefix 自由分布能同时实现所有 Gibbs target；有限容量的共享模型通常只能逼近它们。本章的精确收敛结论适用于前一种模型，扰动定理为后一种模型提供条件性结论，前提是实际输出误差得到测量或上界。
